\documentclass[12pt,a4paper]{article}

% Pakete
\usepackage[utf8]{inputenc}
\usepackage[T1]{fontenc}
\usepackage[ngerman]{babel}
\usepackage{geometry}
\usepackage{minted}        % Für Syntax-Highlighting
\usepackage{xcolor}        % Für Farben

% Seitenränder
\geometry{a4paper, margin=2.5cm}

% Minted-Einstellungen
\setminted{
    frame=single,
    bgcolor=gray!5,
    fontsize=\small,
    linenos=true,          % Zeilennummern
    breaklines=true,       % Lange Zeilen umbrechen
    tabsize=4
}

\title{Mein erstes LaTeX-Dokument mit VHDL}
\author{Dein Name}
\date{\today}

\begin{document}

\maketitle

\section{Einleitung}
Das ist mein erstes Dokument. Hier zeige ich VHDL-Code aus einer externen Datei.

\section{Der VHDL-Code}

% So fügst du die externe VHDL-Datei ein:
% Der Pfad ist relativ zur .tex-Datei!
\inputminted{vhdl}{../rtl/slave.vhd}

\section{Erklärung}
Die Entity \texttt{counter} implementiert einen einfachen Zähler. 
Wichtige Signale sind \mintinline{vhdl}{clk} und \mintinline{vhdl}{reset}.

% Alternativ: Nur bestimmte Zeilen (z.B. Zeile 5-12)
\subsection{Ausschnitt (nur Entity)}
\inputminted[firstline=5, lastline=12]{vhdl}{../rtl/slave.vhd}

\end{document}